\subsection{Reduction of a form with respect to a positive Hermitian form.}

In this section we will assume that E is of finite dimension n over A, and we will denote by Φ a positive nondegenerate Hermitian form on E.

Since the linear maps from E into \(\mathbf{E}^*\) associated with \(\Phi\) are bijective, whatever the elements \(x, y, z\) of E with \(y \neq 0\) , there exists one and only one element \(t\) of E such that \(\Phi(x, z) = \overline{\Phi(t, y)}\) In particular, if \(\mathrm{A} = \mathrm{K}\) or \(\mathrm{A} = \mathrm{K}(i)\) and if \(u\) is a semilinear map for J from E into itself (chap.~II, App. I, no. 1), then for every \(x \in \mathrm{E}\) there exists one and only one element \(u^*(x)\) such that for all \(y \in \mathrm{E}\) , \(\Phi(x, u(y)) = \overline{\Phi(u^*(x), y)}\) We see immediately that \(u^*\) is a semilinear map for J from E into itself; it is called the adjoint of \(u\) .

Remarks. --- 1) When J is the identity, we recover the notion of the adjoint of a homomorphism defined in § 1, no. 8.

\begin{enumerate}
\def\labelenumi{\arabic{enumi})}
\setcounter{enumi}{1}
\tightlist
\item
  Assume always that A = K or A = K(i). Let E \(^{j}\) be the right A-module defined in § 1, no. 2, def. 5. The form Φ is a bilinear form on E × E \(^{j}\) , Φ \(^{j}\) is a bilinear form on E \(^{j}\) × E, and u is an A-linear map from E to E \(^{j}\) . The adjoint of this homomorphism, in the sense of § 1, no. 8, is an A-linear map from E into E \(^{j}\) ; we see at once that this coincides with the map u* defined above.
\end{enumerate}

We say that an endomorphism \(u\) of E is normal (for \(\Phi\) ) if one has \(uu^{*}=u^{*}u\) .
% semantic-fragments: legacy-input-seam
\relax{}
\paragraph{Examples of normal endomorphisms:}

\begin{enumerate}
\def\labelenumi{\arabic{enumi})}
\item
  unitary automorphisms for \(\Phi\) ( \(\S 6\) , no 2), which are characterized by the relation \(u^{-1} = u^{*}\) ( \(\S 1\) , no 8, cor. of prop. 8);
\item
  endomorphisms \(u\) such that \(u^{*}=u\) ; these are called Hermitian endomorphisms.
\end{enumerate}

For every Hermitian endomorphism \(u\) of E, set \(\Phi_{u}(x,y)=\Phi(u(x),y)\) ; we have

\[
\Phi_ {u} (y, x) = \Phi (u (y), x) = \Phi (y, u (x)) = \overline {{\Phi (u (x) , y)}} = \overline {{\Phi_ {u} (x , y)}}.
\]

which shows that \(\Phi_{u}\) is a Hermitian form on E. Conversely, let \(\Psi\) be a Hermitian form on E; since the map \(s_{\Phi}\) from E to E* associated with \(\Phi\) is bijective, there exists, for every \(x \in \mathbf{E}\) , a unique element \(u(x)\) of E such that \(\Psi(x, y) = \Phi(u(x), y)\) ; one verifies easily that \(u\) is a Hermitian endomorphism of E. Thus \(u \to \Phi_{u}\) is a bijection, called canonical, from the set of Hermitian endomorphisms of E onto the set of Hermitian forms on E.

Suppose that A = K or A = K(i); given a bilinear form Ψ on E, there exists, for every \(x \in E\) , a unique element \(u(x)\) , a unique element of E such that Ψ(x, y) = \(\overline{\Phi(u(x), y)}\) ; one sees immediately that u is a semilinear map (for J) from E into itself. Since \(\overline{\Phi(u(x), y)} = \Phi(y, u(x)) = \overline{\Phi(u^{*}(y), x)}\) , one sees that, for Ψ to be symmetric (resp. alternating), it is necessary and sufficient that one have \(u^{*} = u\) (resp. \(u^{*} = -u\) ).

THEOREM 2. --- Let S be a set of endomorphisms of E (resp. of semilinear maps from E into E when A = K or A = K(i)) stable under the map u → u*. Then, if V is a subspace of E stable under S, its orthogonal V⁰ is stable under S. Moreover, E is a direct sum of subspaces stable under S, minimal among the subspaces ≠ \{0\} and stable under S, and pairwise orthogonal.

Indeed, let V be a subspace of E stable under S; for any \(x \in \mathrm{V}^0\) , \(y \in \mathrm{V}\) and \(u \in \mathrm{S}\) , one has \(u^*(y) \in \mathrm{V}\) , whence \(\Phi(y, u(x)) = \Phi(u^*(y), x) = 0\) (resp. \(\Phi(y, u(x)) = \overline{\Phi(u^*(y), x)} = 0\) ), and consequently \(u(x) \in \mathrm{V}^0\) ; thus \(\mathrm{V}^0\) is stable under S, which proves our first assertion. For the second, we proceed by induction on the dimension. \(n\) of E, the case \(n = 0\) being trivial. For \(n \neq 0\) there exists a subspace \(\mathrm{V} \neq \{0\}\) of E stable under S and minimal, for example a stable subspace \(\neq \{0\}\) of minimal dimension. It then suffices to apply to \(\mathrm{V}^0\) the induction hypothesis, since the adjoint of the restriction of \(u\) to \(\mathrm{V}^0\) (with respect to the restriction of \(\Phi\) ) is identical to the restriction to \(\mathrm{V}^0\) of the adjoint of \(u\) .

COROLLARY 1. --- Suppose that A = K or that A = K(i). Let B be a subalgebra of \(\mathfrak{L}_{\mathbf{A}}(\mathbf{E})\) stable under the map \(u \to u^{*}\) . Then E is a semisimple B-module, and is a direct sum of simple submodules that are pairwise orthogonal. The algebra B is semisimple.

Indeed, since every sub-B-module V of E admits a complement, for example \(V^{0}\) the B-module V is semisimple (chap.~VIII, § 3, no. 3, prop. 7). Since every sub-B-module \(\neq\{0\}\) and minimal of E is simple, E is a direct sum of pairwise orthogonal simple submodules. Finally B is a semisimple algebra, since it admits a semisimple and faithful module E whose contramodule is finitely generated (chap.~VIII, § 5, no.~1, prop.~3).

COROLLARY 2. --- With the hypotheses and notations of cor.~1, suppose moreover that the algebra B is commutative. Then all its elements are (normal) semisimple endomorphisms of E. When A = K (resp. A = K(i)), E is a direct sum of pairwise orthogonal simple B-submodules, which are vector spaces of dimension 1 or 2 (resp. 1) over A.

The first assertion follows from chap.~VIII, § 9, no.~1, prop.~2, since B is semisimple. On the other hand every simple B-module is isomorphic to a B-module of the form B/m, where m is a maximal ideal of B, hence to a commutative field L of finite degree over A; when A = K (resp. A = K(i)), the field L is necessarily isomorphic to K or K(i) (resp. K(i)), since K(i) is algebraically closed (chap.~VI, § 2, no.~6, th.~3).

PROPOSITION 4. --- Let u be a normal endomorphism of E. When A is equal to K(i) or to the field of quaternions over K, there exists an orthonormal basis (for Φ) of E consisting of eigenvectors of u. When A = K, u is semisimple and E is a direct sum of subspaces stable under u, pairwise orthogonal, and of dimension 1 or 2.

Let us first examine the case where A is commutative (A = K or A = K(i)). Then the subalgebra B = A{[}u, u*{]} of ℒ\_A(E) is commutative since u is normal; it is stable under the map v → v* by virtue of formulas (32) and (33) of § 1, no.~8. The assertion relative to the case A = K then follows immediately from cor.~2 of th.~2. When A = K(i), this corollary also shows that E is a direct sum of vector subspaces Ax\_i (i = 1, . . ., n) of dimension 1, pairwise orthogonal and stable under u; if we set \(e_{i} = (\Phi(x_{i}, x_{i}))^{-1/2}x_{i}\) , \((e_{i})\) is the desired orthonormal basis.

When A is the field of quaternions over K, it will likewise suffice for us to prove, by virtue of th.~2, that every minimal element of the set of subspaces ≠ \{0\} of E stable under u and u* is of dimension 1. Now such a subspace V necessarily contains an eigenvector x ≠ 0 of u (*), as one sees by noting that the field of quaternions A contains K(i) as an algebraically closed subfield, and by restricting to K(i) the field of scalars of V. It therefore remains only to show that the eigenvector \(x\) of \(u\) is also an eigenvector of \(u^{*}\) . Set \(u(x) = ax (a \in A)\) ; we then have \(\Phi(u(x), x) = a\Phi(x, x) = \Phi(x, x)a = \Phi(x, \bar{a}x)\) since \(\Phi(x, x)\) belongs to the center of A, and on the other hand \(\Phi(u(x), x) = \Phi(x, u^*(x))\) ; it follows that we have \(\Phi(x, u^*(x) - \bar{a}x) = 0\) , and we can therefore write \(u^*(x) = \bar{a}x + z\) , where \(z\) is a vector orthogonal to \(x\) . We therefore have

\[
\Phi (u ^ {*} (x), u ^ {*} (x)) = \bar {a} a \Phi (x, x) + \Phi (z, z) = \Phi (u (x), u (x)) + \Phi (z, z).
\]

Now, since \(u\) is normal, we have

\[
\begin{array}{r l} \Phi (u (x), u (x)) & = \Phi (x, u ^ {*} u (x)) = \Phi (x, u u ^ {*} (x)) \\ & = \Phi (u u ^ {*} (x), x) = \Phi (u ^ {*} (x), u ^ {*} (x)). \end{array}
\]

Consequently we have \(\Phi(z, z) = 0\) , which, by hypothesis, implies \(z = 0\) and \(u^{*}(x) = \bar{a}x\) .Q.E.D.

Remark. --- It follows from th. 2 that the eigenspaces corresponding to two distinct eigenvalues of \(u\) are orthogonal.

PROPOSITION 5. --- Let u be a Hermitian endomorphism of E. The eigenvalues of u belong to K, and there exists an orthonormal basis of E consisting of eigenvectors of u.

When A is equal to K(i) or to the field of quaternions over K, it suffices, by virtue of prop. 4, to prove the first assertion; now, if x is a nonzero eigenvector of u and a is the corresponding eigenvalue, we have, by virtue of the hypothesis u = u*,

\[
a \Phi (x, x) = \Phi (u (x), x) = \Phi (x, u (x)) = \Phi (x, x) \bar {a};
\]

since \(\Phi(x, x)\) is a nonzero element of the center of A, it follows that \(a = \bar{a}\) , hence \(a \in K\) . Consequently a Hermitian matrix has all its eigenvalues in K. Suppose now that A is equal to K. The matrix \(M\) of \(u\) with respect to an orthonormal basis of E is then symmetric; it is therefore a Hermitian matrix if one considers it as a matrix over \(K(i)\) . The first part of the proof then shows that this matrix has all its eigenvalues in K, and hence admits, if \(E \neq \{0\}\) , eigenvectors \(\neq 0\) in E. It follows that every stable subspace for \(u\) which is minimal is necessarily of dimension 1, and the conclusion follows immediately, as in prop. 4.

PROPOSITION 6. --- a) Let Ψ be a Hermitian form (resp. symmetric bilinear form if A = K or A = K(i)) on E. There exists a basis of E which is orthonormal for Φ and orthogonal for Ψ.

\begin{enumerate}
\def\labelenumi{\alph{enumi})}
\setcounter{enumi}{1}
\tightlist
\item
  Suppose A = K or A = K(i), and let Ψ be an alternating bilinear form on E. There exists a basis of E which is orthonormal for Φ and with respect to which the matrix of Ψ is of the form
\end{enumerate}

\[
\left( \begin{array}{c c c c} 0 & a _ {1} & 0 & 0 \ldots 0 \\ - a _ {1} & 0 & 0 & 0 \ldots 0 \\ 0 & 0 & 0 & a _ {2} \ldots 0 \\ 0 & 0 & - a _ {2} & 0 \ldots 0 \\ \ldots \ldots \end{array} \right).
\]

where the \(a_{i}\) are \(\geqslant 0\) in K.

When \(\Psi\) is Hermitian, our assertion follows immediately from prop. 5 and the canonical correspondence between Hermitian forms on E and Hermitian endomorphisms for \(\Phi\) . In the other two cases, let \(u\) be the semilinear map from E to itself defined at the beginning of this n° by the formula \(\Psi(x, y) = \overline{\Phi(u(x), y)}\) . Then \(u^2\) is an A-linear map; when \(\Psi\) is symmetric (resp. alternating), one has \(u = u^*\) (resp. \(u = -u^*\) ), hence \(u^2\) is Hermitian; whence also

\[
\begin{array}{r l} \Phi (u ^ {2} (x), x) & = \overline {{\Phi (x , u ^ {2} (x))}} = \Phi (u ^ {*} (x), u (x)) = \Phi (u (x), u (x)) \\ & (\text { resp. } - \Phi (u (x), u (x))) \end{array}
\]

this shows that the Hermitian form \((x, y) \to \Phi(u^{2}(x), y)\) is positive (resp. negative). Then apply th. 2, and let V be a minimal element of the family of subspaces \(\neq\{0\}\) of E stable under u and under \(u^{*}\) . Since \(u^{2}\) is a Hermitian endomorphism, V contains an eigenvector \(x \neq 0\) of \(u^{2}\) ; set \(u^{2}(x) = ax\) , where \(a \in K\) (prop. 5).

When \(\Psi\) is symmetric, the inequality \(\Phi(u^{2}(x), x) \geqslant 0\) shows that one has \(a \geqslant 0\) . Set \(y = a^{1/2}x + u(x)\) ; we have \(u(y) = a^{1/2}u(x) + ax = a^{1/2}y\) , and Ay is stable under \(u\) and \(u^{*} = u\) . If \(y = 0\) , Ax is stable under \(u\) and \(u^{*}\) ; in all cases, V is a subspace of dimension 1, and our conclusion follows immediately from th. 2.

When \(\Psi\) is alternating, one has \(a \leqslant 0\) ; set

\[
y = (- a) ^ {1 / 2} x + u (x), \quad z = (- a) ^ {1 / 2} x - u (x).
\]

\[
\text { One has } u (y) = - (- a) ^ {1 / 2} z, u (z) = (- a) ^ {1 / 2} y \text { and }
\]

\[
\Phi (y, z) = - a \Phi (x, x) - \Phi (u (x), u (x)) = - \Phi (a x, x) + \Phi (u ^ {2} (x), x) = 0.
\]

If \(y = z = 0\) , one has \(a = 0\) and \(u(x) = 0\) , hence \(Ax\) is stable under \(u\) and \(u^{*}\) , \(V\) is of dimension 1, and the matrix of the restriction of \(\Psi\) to \(V\) is zero since \(\Psi\) is alternating. Otherwise, \(y\) and \(z\) are both \(\neq 0\) , they generate \(V\) , and since they are orthogonal, \(V\) is of dimension 2 and the matrix of the restriction of \(\Psi\) to \(V\) is of the form \(\begin{pmatrix} 0 & b \\ -b & 0 \end{pmatrix}\) . Finally, by virtue of the formula \(\Psi(x', y') = \overline{\Phi(u(x'), y')}\) , one has \(\Psi(x', y') = 0\) when \(x'\) and \(y'\) belong to two subspaces stable under \(u\) and orthogonal for \(\Phi\) . This proves the existence of an orthonormal basis of \(E\) (for \(\Phi\) ) with respect to which the matrix of \(\Psi\) has the indicated form. Q.E.D.

Exercises. --- 1) Suppose the hypotheses at the beginning of § 7 are satisfied; let \(\Phi\) be a positive Hermitian form on E.

\begin{enumerate}
\def\labelenumi{\alph{enumi})}
\item
  For us to have \(\Phi(x,x)\Phi(y,y)=\Phi(x,y)\overline{\Phi(x,y)}\) , it is necessary and sufficient that \(x\) and \(y\) be linearly dependent, or that the plane spanned by \(x\) and \(y\) be isotropic.
\item
  We assume that, for all \(x \in E\) , \(\Phi(x, x)\) is a square in K. Show that for any two vectors x, y in E, we have
\end{enumerate}

\[
\sqrt {\Phi (x + y , x + y)} \leqslant \sqrt {\Phi (x , x)} + \sqrt {\Phi (y , y)}.
\]

If \(\Phi\) is nondegenerate, the two sides of this inequality can be equal only if \(\alpha x + \beta y = 0\) , where \(\alpha\) and \(\beta\) are two elements of K, not both zero, and such that \(\alpha \beta \leqslant 0\) .

\begin{enumerate}
\def\labelenumi{\arabic{enumi})}
\setcounter{enumi}{1}
\tightlist
\item
  We assume \(\mathbf{A} = \mathbf{K}\) or \(\mathbf{A} = \mathbf{K}(i)\) . Let \(X\) be a square Hermitian matrix of order \(n\) on \(\mathbf{A}\) such that for every nonempty subset \(\mathbf{H}\) of \([1, n]\) one has \(X_{\mathbf{H},\mathbf{H}} \geqslant 0\) (notations of prop. 3 of no. 1).
\end{enumerate}

\begin{enumerate}
\def\labelenumi{\alph{enumi})}
\item
  Let \(\lambda\) be an element \(>0\) of K; show that the matrix \(X + \lambda I\) is positive nondegenerate (use prop. 3 of no. 1, arguing by induction on \(n\) ).
\item
  Deduce that the Hermitian matrix X is positive.
\end{enumerate}

\begin{enumerate}
\def\labelenumi{\arabic{enumi})}
\setcounter{enumi}{2}
\item
  We assume that \(\mathbf{A} = \mathbf{K}\) or \(\mathbf{A} = \mathbf{K}(i)\) and that E is finite-dimensional. Let \(\Phi_1, \Phi_2\) be two positive Hermitian forms on E, and let \(V = (\alpha_{ij})\) , \(W = (\beta_{ij})\) be the matrices of these forms with respect to the same basis \((e_i)\) of E. Show that the Hermitian form \(\Phi\) whose matrix \((\gamma_{ij})\) with respect to \((e_i)\) is such that \(\gamma_{ij} = \alpha_{ij}\beta_{ij}\) for every pair of indices, is positive; moreover, if \(\Phi_1\) and \(\Phi_2\) are nondegenerate, the same holds for \(\Phi\) . (In the computation of \(\Phi(x, x)\) express the \(\alpha_{ij}\) in terms of the values \(\Phi_1(c_i, c_i)\) for an orthogonal basis \((c_i)\) of E relative to \(\Phi_1\) .)
\item
  Suppose that A = K or A = K(i). Let R be a Hermitian matrix of order n over A; we say that R has signature \((s, t)\) if the Hermitian form on \(A^{n}\) having R as its matrix with respect to the canonical basis, has signature \((s, t)\) . We denote by \(\Delta_{k}\) the principal minor of R obtained by deleting from R the rows and columns with index \textgreater{} k ; we assume that \(\Delta_{s+t} \neq 0\) and that, for no index \(k < s + t\) , \(\Delta_{k}\) and \(\Delta_{k+1}\) are not simultaneously zero (cf.~§ 6, exerc. 1 b)). Show that if \(\Delta_{k} = 0\) for a \(k < s + t\) , \(\Delta_{k-1}\) and \(\Delta_{k+1}\) have opposite signs (method of exerc. 1 a) of § 6) and that the number s - t is equal to
\end{enumerate}

\[
\operatorname{sgn} \Delta_ {1} + \operatorname{sgn} \left(\Delta_ {1} \Delta_ {2}\right) + \dots + \operatorname{sgn} \left(\Delta_ {s + t - 1} \Delta_ {s + t}\right)
\]

(use exercise 2 of § 6).

\begin{enumerate}
\def\labelenumi{\arabic{enumi})}
\setcounter{enumi}{4}
\item
  We assume that A = K or A = K(i); let R, S be two Hermitian matrices over A, with signatures \((s, t)\) and \((s', t')\) respectively (exerc. 4); show that the matrix \(R \otimes S\) is a Hermitian matrix of signature \((ss' + tt', st' + s't)\) .
\item
  Let K be a maximal ordered field, L a finite-dimensional simple algebra over K. If \((s, t)\) is the signature of the symmetric bilinear form \((x, y) \to \mathrm{Tr}_{\mathrm{L/K}}(xy)\) on L, show that we have:
\end{enumerate}

\(s-t=m\) if L is isomorphic to a matrix algebra of order \(m\) over K;

s - t = 0 if L is isomorphic to a matrix algebra over the field K(i);

\(s - t = -2m\) if L is isomorphic to a matrix algebra of order \(m\) over the field of quaternions over K.

\begin{enumerate}
\def\labelenumi{\arabic{enumi})}
\setcounter{enumi}{6}
\tightlist
\item
  Assume that A satisfies the conditions at the beginning of §7 and that E has finite dimension n. Let Φ be a positive nondegenerate Hermitian form on E.
\end{enumerate}

\begin{enumerate}
\def\labelenumi{\alph{enumi})}
\item
  Show that every similarity for \(\Phi\) is written in a unique way as the product of a homothety with ratio \(\geqslant 0\) in K, and of a unitary transformation.
\item
  For every basis \((a_{i})_{1\leqslant i\leqslant n}\) of E, show that there exists one and only one orthonormal basis \((e_{i})_{1\leqslant i\leqslant n}\) of E satisfying the following conditions: \(1^{\circ}\) for all \(m\) such that \(1\leqslant m\leqslant n\) , the subspace spanned by \(a_1,\ldots ,a_m\) is identical to the subspace spanned by \(e_1,\ldots ,e_m\) ; \(2^{\circ}\) we have \(\Phi (a_i,e_i) > 0\) in K, for every index \(i\) (cf. § 6, no. 1, prop. 1).
\item
  Deduce from \(b\) that for every invertible square matrix \(M\) of order \(n\) on A, there exists one and only one pair of matrices \((L, U)\) of order \(n\) such that \(U\) is a unitary matrix, that \(L = (\lambda_{ij})\) has only zeros below its diagonal and diagonal entries \(\lambda_{ii}\) belonging to K and \(>0\) and finally that \(M = LU\) .
\end{enumerate}

Assume A = K; let L be a finite-dimensional Hermitian space over Λ, whose metric form is positive non-degenerate. Let M, N be two subsets of L, and u a bijection from M onto N such that one has \(e(u(a), u(b)) = e(a, b)\) for any points a, b in M, show that there exists a displacement whose restriction to M is equal to u (argue by induction on the dimension of the affine linear variety spanned by M, as in the Gram-Schmidt orthogonalization process). Does the proposition extend to the case where A equals K(i) or the field of quaternions over K?

\begin{enumerate}
\def\labelenumi{\arabic{enumi})}
\setcounter{enumi}{8}
\item
  Assume the conditions of No. 3 are satisfied, and in addition that A is the field of quaternions over K. Let u be a normal endomorphism of E. Show that E is the direct sum of subspaces \(F_{k}\) ( \(1 \leqslant k \leqslant r\) ), pairwise orthogonal, such that in each of the \(F_{k}\) there exists an orthonormal basis ( \(e_{ik}\) ) ( \(1 \leqslant i \leqslant n_{k}\) ) and that one has \(u(e_{ik}) = \lambda_{k}e_{ik}\) for \(1 \leqslant i \leqslant n_{k}\) , two \(\lambda_{k}\) with distinct indices not being transformed into each other by an inner automorphism of A. Moreover, if ( \(F_{k}'\) ) is a second decomposition of E having the same properties, with a system ( \(\lambda_{k}'\) ) of eigenvalues, one has \(F_{k}' = F_{k}\) (up to a permutation of the indices) and \(\lambda_{k}' = \alpha_{k}\lambda_{k}\alpha_{k}^{-1}\) ; the set of eigenvectors of u for the eigenvalue \(\lambda_{k}\) is the subspace over the commutant \(A_{k}\) of \(\lambda_{k}\) in A, generated by the \(e_{ik}\) ( \(1 \leqslant i \leqslant n_{k}\) ).
\item
  Assume the conditions of No. 3 are satisfied, and in addition that A is equal to K(i) or to the field of quaternions over K. Let u be a normal endomorphism of E.
\end{enumerate}

\begin{enumerate}
\def\labelenumi{\alph{enumi})}
\item
  Show that for a vector subspace F of E to be such that \(u(\mathrm{F}) \subset \mathrm{F}\) , it is necessary and sufficient that F be generated by eigenvectors of u; we then have \(u(\mathrm{F}^{0}) \subset \mathrm{F}^{0}\) and \(u^{*}(\mathrm{F}) \subset \mathrm{F}\) .
\item
  For \(u\) to be unitary (resp. for \(u^{*} = u\) , \(u^{*} = -u\) ), it is necessary and sufficient that for every eigenvalue \(\lambda\) of \(u\) one has \(\lambda \bar{\lambda} = 1\) (resp. \(\bar{\lambda} = \lambda\) , \(\bar{\lambda} = -\lambda\) ).
\item
  We say that a Hermitian endomorphism \(u\) of E is positive (resp. positive and nondegenerate) if the Hermitian form \((x, y) \to \Phi(u(x), y)\) canonically corresponding to it is positive (resp. positive and nondegenerate); for this it is necessary and sufficient that all eigenvalues of \(u\) be \(\geqslant 0\) (resp. \(>0\) ).
\item
  Show that for every integer m\textgreater{} 0, there exists a normal endomorphism v of E such that \(v^{m}=u\) . If u is positive Hermitian, there exists a unique positive Hermitian endomorphism v such that \(v^{m}=u\) , and there exists a polynomial \(f\in K[X]\) such that \(v=f(u)\) ; this last result is also valid when A = K.
\end{enumerate}

\begin{enumerate}
\def\labelenumi{\arabic{enumi})}
\setcounter{enumi}{10}
\tightlist
\item
  Assume the conditions of no. 3 are satisfied, and in addition that A = K. Let u be a normal endomorphism of E.
\end{enumerate}

\begin{enumerate}
\def\labelenumi{\alph{enumi})}
\item
  Let V be a minimal element of the set of subspaces of E not reduced to 0, stable under u and \(u^{*}\) ; show that if V has dimension 2, the restriction of u to V is a direct similarity with multiplier \textgreater{} 0.
\item
  Show that every subspace of E stable under u is also stable under \(u^{*}\) . (Let \(E_{0}\) be the vector space obtained from E by extending the field of scalars to \(\mathrm{K}(i)\) ; \(\Phi\) is the restriction to E of a positive nondegenerate Hermitian form on \(E_{0}\) and u the restriction to E of a normal endomorphism \(u_{0}\) of \(E_{0}\) (for this form); moreover, E is the set of \(x \in E_{0}\) invariants under an involutive semilinear bijection j of \(E_{0}\) , and we have \(u_{0}j = ju_{0}\) . Then apply exerc. 10 a).)
\end{enumerate}

\begin{enumerate}
\def\labelenumi{\arabic{enumi})}
\setcounter{enumi}{11}
\item
  Assume the conditions of No. 3 are satisfied, and in addition that A is equal to K(i) or to the field of quaternions over K. Show that for every endomorphism u of E, there exists an orthonormal basis of E with respect to which the matrix of u has only zeros below the diagonal. (Proceed by induction on the dimension of E, considering an eigenvector of u.)
\item
  Let A be a field satisfying the conditions at the beginning of § 7, and let E, F be two finite-dimensional vector spaces over A, \(\Phi\) (resp. \(\Psi\) ) a positive non-degenerate Hermitian form on E (resp. F). Show that, for every linear map \(u\) from E to F, the adjoint \(u^{*}\) of \(u\) (for \(\Phi\) and \(\Psi\) ; cf. § 1, no. 8) is such that \(u^{*}u\) and \(uu^{*}\) are positive Hermitian endomorphisms (exercise 10 c)) of E and F, respectively.
\item
  Assume the conditions of No. 3 are satisfied. Let u be an endomorphism of E, \(h_{1}\) , \(h_{2}\) the two positive Hermitian endomorphisms of E, such that we have \(h_{1}^{2}=u^{*}u\) , \(h_{2}^{2}=uu^{*}\) (Exercises 13 and 10 d)).
\end{enumerate}

\begin{enumerate}
\def\labelenumi{\alph{enumi})}
\item
  Show that there exists a unitary endomorphism \(\nu\) such that \(u = \nu h_{1} = h_{2}\nu\) and in particular that \(h_{1}\) and \(h_{2}\) are similar (note that \(\bar{u}(0) = \bar{h}_{1}^{-1}(0)\) and that if V is the subspace orthogonal to \(\bar{u}(0)\) , one has \(\Phi(u(x), u(x)) = \Phi(h_{1}(x), h_{1}(x))\) for all \(x \in \mathrm{V}\) . For \(\nu\) to be uniquely determined, it is necessary and sufficient that \(u\) be bijective. In order for one to be able to take \(\nu\) commuting with \(h_{1}\) , it is necessary and sufficient that \(u\) be normal.
\item
  From a), deduce that every square matrix M of order n over A can be written as UDV, where U and V are unitary matrices and D is a diagonal matrix whose entries are \(\geqslant 0\) and have as their squares the eigenvalues of MM*.
\end{enumerate}

\begin{enumerate}
\def\labelenumi{\arabic{enumi})}
\setcounter{enumi}{14}
\tightlist
\item
  Assume the conditions of No. 3 are satisfied. Show that every positive Hermitian matrix H over A can be written in the form \(LL^{*}\) , where \(L = (\lambda_{ij})\) has only zeros below its diagonal and diagonal terms belonging to K and \(\geqslant 0\) ; moreover, L is uniquely determined by these conditions when H is invertible (cf.~exerc. 10 d) and 7 c).
\end{enumerate}

¶ 16) We assume that the conditions of no. 3 are satisfied and, in addition, that A = K(i). Let u be an endomorphism of E.

\begin{enumerate}
\def\labelenumi{\alph{enumi})}
\item
  The set of eigenvalues of \(u\) is contained in the set U of values of \(\Phi(x, u(x))\) when \(x\) ranges over the set of elements of E such that \(\Phi(x, x) = 1\) .
\item
  A subset C of A = K(i) is said to be convex if, for every pair of elements \((\xi, \eta) \in \mathrm{C}^{2}\) and every \(\tau \in K\) such that \(0 \leqslant \tau \leqslant 1\) we have \(\tau\xi + (1 - \tau)\eta \in C\) . Show that the set U is convex. (Reduce to the case where n = 2; by writing u in the form \(\nu + iw\) , where \(\nu\) and w are Hermitian, and replacing u if necessary by \(\lambda u\) , where \(\lambda \in A\) and \(\lambda\overline{\lambda} = 1\) show that everything reduces to proving the following property. Let \(f(\xi_{1}, \xi_{2}) = \xi_{1}\overline{\xi}_{1} + \xi_{2}\overline{\xi}_{2}\) , \(g(\xi_{1}, \xi_{2}) = a\xi_{1}\overline{\xi}_{1} + b\xi_{2}\overline{\xi}_{2}\) ( \(a \in K, b \in K\) ), \(h(\xi_{1}, \xi_{2}) = \alpha\xi_{1}\overline{\xi}_{1} + \beta\xi_{1}\overline{\xi}_{2} + \beta\xi_{2}\overline{\xi}_{1} + \gamma\xi_{2}\overline{\xi}_{2}\) ( \(\alpha \in K, \gamma \in K, \beta \in A\) ); let \((\eta_{1}, \eta_{2}) \in A^{2}\) , \((\zeta_{1}, \zeta_{2}) \in A^{2}\) such that \(f(\eta_{1}, \eta_{2}) = f(\zeta_{1}, \zeta_{2}) = 1\) , \(g(\eta_{1}, \eta_{2}) = g(\zeta_{1}, \zeta_{2}) = 1\) , \(h(\eta_{1}, \eta_{2}) > 0\) , \(h(\zeta_{1}, \zeta_{2}) < 0\) ; there then exists \((\theta_{1}, \theta_{2}) \in A^{2}\) such that \(f(\theta_{1}, \theta_{2}) = 1\) , \(g(\theta_{1}, \theta_{2}) = 1\) and \(h(\theta_{1}, \theta_{2}) = 0\) . One will begin by noting that for every pair \((\xi_{1}, \xi_{2})\) , there exist \(\mu \in A\) such that \(\mu\overline{\mu} = 1\) and \(\beta\mu\xi_{1}\overline{\xi}_{2} + \beta\overline{\mu}\xi_{2}\overline{\xi}_{1} = 0\) , which will make it possible to reduce to the case where \(\beta = 0\) ; use prop. 5 of chap.~VI, § 2, no.~5.)
\item
  Show that if \(u\) is normal, U is the smallest convex set containing all the eigenvalues of \(u\) . Give an example where \(u\) is not normal but where U still possesses the preceding property. (Take as eigenvalues of \(u\) the elements \(\pm 1 \pm i\) as simple eigenvalues, and 0 as a double eigenvalue.)
\end{enumerate}

¶ 17) We assume the conditions of no.~3 are satisfied; for every \(\xi\in\mathbf{A}\) , we denote by \(|\xi|\) the element \(\rho\geqslant0\) of K such that \(\rho^{2}=\xi\bar{\xi}\) (absolute value of \(\xi\) ). For every square matrix \(M=(\alpha_{ij})\) over A, we set \(f(M)=\max_{i}|\alpha_{ii}|\) , \(g(M)=\max_{i,j}|\alpha_{ij}|\) , and denote by \(\varphi(M)\) the largest absolute value of the eigenvalues of \(M\) (it will be shown that this definition makes sense when A is the field of quaternions over K).

\begin{enumerate}
\def\labelenumi{\alph{enumi})}
\tightlist
\item
  Let \(A, B, D\) be three square matrices of order \(n\) over A. Show that if \(D\) is diagonal, then
\end{enumerate}

\[
g ^ {2} (A D B ^ {*}) \leqslant f ^ {2} (D) f (A A ^ {*}) f (B B ^ {*})\tag{1}
\]

(use inequality (1) of no. 1). Deduce that

\[
g (A B B ^ {*} A ^ {*}) \leqslant f (A A ^ {*}) \varphi (B B ^ {*})\tag{2}
\]

(apply prop. 5 to the Hermitian matrix \(BB^{*}\) ). Deduce that, for \(m\) arbitrary square matrices \(A_{i}\) ( \(1 \leqslant i \leqslant m\) ) over A, we have

\begin{enumerate}
\def\labelenumi{(\arabic{enumi})}
\setcounter{enumi}{2}
\tightlist
\item
\end{enumerate}

\[
g ^ {2} (A _ {1} A _ {2} \dots A _ {m}) \leqslant f (A _ {1} A _ {1} ^ {*}) \varphi (A _ {2} A _ {2} ^ {*}) \dots \varphi (A _ {m - 1} A _ {m - 1} ^ {*}) f (A _ {m} ^ {*} A _ {m})\tag{4}
\]

\[
\varphi^ {2} (A _ {1} A _ {2} \dots A _ {m}) \leqslant \varphi (A _ {1} A _ {1} ^ {*}) \varphi (A _ {2} A _ {2} ^ {*}) \dots \varphi (A _ {m} A _ {m} ^ {*})\tag{5}
\]

\[
g ^ {2} (A _ {1} A _ {2} \dots A _ {m}) \leqslant \varphi (A _ {1} A _ {1} ^ {*}) \dots \varphi (A _ {m} A _ {m} ^ {*}).
\]

(For (3), reason by induction starting from (2). For (4), use exerc. 12; finally deduce (5) from (3).)

Show that inequality (3) no longer necessarily holds when one replaces \(A_{m}^{*}A_{m}\) by \(A_{m}A_{m}^{*}\) (observe that one can have \(f(A^{*}A)\neq f(A A^{*})\) ), or when one replaces \(\varphi(A_{i}A_{i}^{*})\) by \(f(A_{i}A_{i}^{*})\) for \(2\leqslant i\leqslant m-1\) (take all the \(A_{i}\) equal to the square matrix all of whose elements are 1).

\begin{enumerate}
\def\labelenumi{\alph{enumi})}
\setcounter{enumi}{1}
\tightlist
\item
  Let u be a normal endomorphism of E; if λ is an eigenvalue of u (an element of K(i) when A = K), then \(\lambda\bar{\lambda}\) is an eigenvalue of u*u. For every normal matrix M (matrix of a normal endomorphism of E with respect to an orthonormal basis), we therefore have φ²(M) = φ(MM*). Deduce that one also has g(M) ≤ φ(M), and more generally, if M₁, . . ., Mₘ are normal,
\end{enumerate}

\begin{enumerate}
\def\labelenumi{(\arabic{enumi})}
\setcounter{enumi}{5}
\tightlist
\item
\end{enumerate}

\[
g (M _ {1} \dots M _ {m}) \leqslant \varphi (M _ {1}) \dots \varphi (M _ {m})\tag{7}
\]

\[
\varphi (M _ {1} \dots M _ {m}) \leqslant \varphi (M _ {1}) \dots \varphi (M _ {m})
\]

(use (4) and (5)).

\begin{enumerate}
\def\labelenumi{\alph{enumi})}
\setcounter{enumi}{2}
\tightlist
\item
  Show that if \(H\) is a positive Hermitian matrix over A, we have \(f(H) = g(H) \leqslant \varphi(H)\) (use (3) and b), by writing \(H = AA^{*}\) ).
\end{enumerate}

¶ 18) We assume the conditions of no. 3 are satisfied.

\begin{enumerate}
\def\labelenumi{\alph{enumi})}
\tightlist
\item
  For every endomorphism u of E, let \((e_{i})\) be an orthonormal basis of E consisting of eigenvectors of \(u^{*}u\) , and let \(\rho_{i}\) be the eigenvalue of \(u^{*}u\) corresponding to the vector \(e_{i}\) ; we set \(s(u) = (\sum_{i=1}^{n} \rho_{i})^{1/2}\) (square root of \(\mathrm{Tr}(u^{*}u)\) when A is commutative); we have \(s(u^{*}) = s(u)\) . If u, v are two endomorphisms of E, and U, V their matrices with respect to the same orthonormal basis of E, show that for the matrix \(UV^{*} + VU^{*}\) , with elements in K, one has \((\mathrm{Tr}(UV^{*} + VU^{*}))^{2} \leqslant 4s(u)s(v)\) (note that \((u^{*} + \lambda v^{*})(u + \lambda v)\) is positive Hermitian for every \(\lambda \in K\) ); deduce that \(s(u + v) \leqslant s(u) + s(v)\) . If in addition A is commutative, one has
\end{enumerate}

\[
\left| \operatorname{Tr} (u v) \right| ^ {2} \leqslant s (u) s (v) \quad \text { and } \quad \left| \operatorname{Tr} (u) \right| \leqslant \sqrt {n} s (u)
\]

(write \(u\) as a product using Exerc. 14 b)).

\begin{enumerate}
\def\labelenumi{\alph{enumi})}
\setcounter{enumi}{1}
\tightlist
\item
  We further assume A is commutative (hence equal to K or K(i)). If \(H_{1}, H_{2}\) are two positive Hermitian square matrices, show that one has \(\mathrm{Tr}(H_{1}H_{2}) \geqslant 0\) (reduce to the case where \(H_{2}\) is a diagonal matrix). Deduce that we have (using the notation of exercise 17)
\end{enumerate}

\[
\left| \operatorname{Tr} (H _ {1} H _ {2}) \right| \leqslant \varphi (H _ {1}) \operatorname{Tr} (H _ {2}) \leqslant \operatorname{Tr} (H _ {1}) \operatorname{Tr} (H _ {2}).
\]

From these inequalities, conclude that for any two endomorphisms \(u, \nu\) of E, we then have \(s(u\nu) \leqslant s(u)s(\nu)\) .

\begin{enumerate}
\def\labelenumi{\alph{enumi})}
\setcounter{enumi}{2}
\tightlist
\item
  Assuming A is still commutative, let \(\prod_{i=1}^{n}(x-\lambda_i)\) be the decomposition into linear factors of the characteristic polynomial of an arbitrary endomorphism \(u\) of E; show that one has \(\sum_{i=1}^{n}|\lambda_i|^2\leqslant(s(u))^2\) and that, for the two sides of this inequality to be equal, it is necessary and sufficient that \(u\) is normal (use Exerc. 12).
\end{enumerate}

\begin{enumerate}
\def\labelenumi{\arabic{enumi})}
\setcounter{enumi}{18}
\item
  Assume the conditions of no. 3 are satisfied. Let M be a square matrix of order n over A; show that for every square submatrix N of M (obtained by deleting from M a certain number of rows and the columns with the same indices as those rows), one has (with the notations of exercise 17) \(\varphi^{2}(N) \leqslant \varphi(MM^{*})\) (apply formula (4) of exerc. 17 appropriately). If in particular M is a normal matrix, \(\varphi(N) \leqslant \varphi(M)\) (cf. exercise 17 b)).
\item
  We assume that the conditions of No. 3 are satisfied and, in addition, that A is equal to K or to K(i). Apply the results of Exercises 17 to 19 to the extension of Φ to the p-th exterior powers (§ 1, No. 9) and to the p-th exterior powers of the endomorphisms or matrices under consideration. In particular, show that, if \(\prod_{i=1}^{n}(X-\lambda_i)\) and \(\prod_{i=1}^{n}(X-\rho_i^2)\) are the decompositions into linear factors of the characteristic polynomials of an
\end{enumerate}

endomorphism \(u\) of E and of the endomorphism \(u^*u\) and if we suppose \(|\lambda_i| \geqslant |\lambda_{i+1}|\) and \(\rho_i \geqslant \rho_{i+1} \geqslant 0\) for \(1 \leqslant i \leqslant n-1\) , one has

\[
\mid \lambda_ {1} \lambda_ {2} \dots \lambda_ {h} \mid \leqslant \rho_ {1} \rho_ {2} \dots \rho_ {h}
\]

for \(1 \leqslant h \leqslant n - 1\) and \(|\lambda_1\lambda_2 \ldots \lambda_n| = \rho_1\rho_2 \ldots \rho_n\) .

\begin{enumerate}
\def\labelenumi{\arabic{enumi})}
\setcounter{enumi}{20}
\tightlist
\item
  We suppose the conditions of No. 3 are satisfied. Let u be a Hermitian endomorphism of E and let \(\prod_{i=1}^{n}(X-\lambda_{i})\) be the decomposition into linear factors of its characteristic polynomial; we assume \(\lambda_{i}\geqslant\lambda_{i+1}\) for \(1\leqslant i\leqslant n-1\) .
\end{enumerate}

\begin{enumerate}
\def\labelenumi{\alph{enumi})}
\item
  Show that the largest (resp. the smallest) of the eigenvalues \(\lambda_{i}\) in K is equal to the largest (resp. the smallest) of the values of \(\Phi(u(x), x)\) when \(x\) ranges over the set of \(x \in E\) such that \(\Phi(x, x) = 1\) . (Reason directly, or apply Exerc. 16 c) by reducing to the case where A = K(i) ( \(\S\) 3, exercise 4).)
\item
  Let \(\Psi_{\mathrm{v}}\) be the restriction to a vector subspace V of E of the Hermitian form \(\Psi\) associated with \(u\) , and let \(u_{\mathrm{v}}\) be the Hermitian endomorphism of V associated with \(\Psi_{\mathrm{v}}\) . Show that \(\lambda_{k}\) is the smallest of the largest eigenvalues of the \(u_{\mathrm{v}}\) as V ranges over the set of vector subspaces of E of dimension \(n - k + 1\) (using Prop. 5).
\end{enumerate}

\begin{enumerate}
\def\labelenumi{\arabic{enumi})}
\setcounter{enumi}{21}
\tightlist
\item
  We assume that the conditions of no. 3 are satisfied, and in addition that A is equal to K(i) or to the field of quaternions over K. Let u, v be two normal endomorphisms of E; let (E \(_{i}\) ) \(_{1\leqslant i\leqslant r}\) (resp. (F \(_{j}\) ) \(_{1\leqslant j\leqslant s}\) ) be the decomposition of E into a direct sum of pairwise orthogonal subspaces, such that in E \(_{i}\) (resp. F \(_{j}\) ) there is an orthonormal basis formed by eigenvectors of u (resp. v) for the same eigenvalue \(\lambda_{i}\) (resp. \(\mu_{j}\) ), and that \(\lambda_{h}\) and \(\lambda_{i}\) (resp. \(\mu_{j}\) and \(\mu_{k}\) ) are not transformed into one another by an inner automorphism of A if \(h\neq i\) (resp. \(j\neq k\) ) (cf. prop. 4 and exerc. 9). For an endomorphism w of E to satisfy uw = wν, it is necessary and sufficient that for every j ( \(1\leqslant j\leqslant s\) ), w(F \(_{j}\) ) be contained in one of the E \(_{i}\) and that the image under w of every eigenvector of v relative to the eigenvalue \(\mu_{j}\) be an eigenvector of u relative to the eigenvalue \(\mu_{j}\) (which implies in particular that \(\mu_{j}\) and \(\lambda_{i}\) are transformed into one another by an inner automorphism of A). Deduce that if this is so, then one has u*w = wν*.
\end{enumerate}

¶ 23) We assume that the conditions of no. 3 are satisfied. Let u and v be two endomorphisms of E.

\begin{enumerate}
\def\labelenumi{\alph{enumi})}
\item
  We assume that u and uv are normal. For vu to be normal, it is necessary and sufficient that v and \(u^{*}u\) commute. (To see that the condition is necessary, use the relation \(u(vu) = (uv)u\) and exerc. 22; to see that it is sufficient, use exerc. 14 a) and 10 d).
\item
  We assume \(u, v\) and \(uv\) normal; let \(\beta\) be the largest eigenvalue of \(uu^*\) , and let \(F\) be the subspace of \(E\) formed by the eigenvectors of \(uu^*\) relative to the eigenvalue \(\beta\) . Show that \(v(F) \subset F\) . (Note that for every normal endomorphism \(u\) and every \(x \in E\) , one has
\end{enumerate}

\[
\Phi (u (x), u (x)) = \Phi (u ^ {*} (x), u ^ {*} (x)).
\]

Reduce to the case where A is equal to K(i) or to the field of quaternions over K; F then admits a basis formed by eigenvectors for u (and \(u^{*}\) ); note that for such a vector z, we have \(\Phi(u^{*}u\nu(z),\nu(z))=\beta\Phi(\nu(z),\nu(z))\) Deduce from this that \(\nu u\) is normal (argue by induction on the number of distinct eigenvalues of \(u^{*}u\) ). If h (resp. \(h'\) ) is the positive Hermitian endomorphism such that \(h^{2}=uu^{*}\) (resp. \(h'^{2}=\nu\nu^{*}\) ) and if we set \(u=hu_{1}\) , \(\nu=h'\nu_{1}\) , where \(u_{1}\) and \(\nu_{1}\) are unitary, h commuting with \(u_{1}\) and \(h'\) with \(\nu_{1}\) (exerc. 14 a)), show that the pairs \((h,h')\) , \((h,\nu_{1})\) and \((h',u_{1})\) are permutable; and conversely. Deduce that \(u^{m}\nu^{n}\) , \(\nu^{n}u^{m}\) , \(u\nu^{*}\) and \(\nu^{*}u\) are then normal (m and n integers \textgreater0 arbitrary).

¶ 24) Assume the conditions of no. 3 are satisfied. Let Γ be a group of automorphisms of E such that every \(u \in \Gamma\) is normal. Show that there exists a decomposition of E into a direct sum of subspaces \(E_k\) ( \(1 \leqslant k \leqslant r\) ) pairwise orthogonal, such that the restriction to \(E_k\) of every \(u \in \Gamma\) is of the form \(\lambda_k v_k\) , where \(\lambda_k\) is an element \textgreater0 of K and where \(v_k\) is a unitary endomorphism of \(E_k\) (decompose each \(u \in \Gamma\) in the form \(h\nu\) , where \(v\) is unitary, \(h\) positive Hermitian, and \(h^2 = uu^*\) (Exerc. 14a); use Exerc. 23b and apply Cor. 2 of Thm. 2 to the algebra generated by the Hermitian endomorphisms \(h\) corresponding to the \(u \in \Gamma\) ). Deduce that if Γ is finite, the elements \(u \in \Gamma\) are unitary.

\begin{enumerate}
\def\labelenumi{\arabic{enumi})}
\setcounter{enumi}{24}
\tightlist
\item
  Suppose that A = K (maximal ordered field). Let L be an n-dimensional Euclidean space over A, whose metric form is positive non-degenerate.
\end{enumerate}

Let S be a nondegenerate affine quadric in L (§ 6, exerc. 25). If S admits a center a and if one takes a as origin in L, show that there exists an orthonormal basis \((e_{i})\) for L such that, with respect to this basis, S has the equation \(\lambda_{1}\xi_{1}^{2} + \cdots + \lambda_{n}\xi_{n}^{2} = 1\) . Moreover, if two orthonormal bases have this property, the elements \(\lambda_{i} \in K\) which correspond to them are the same up to order.

If S admits no center, show that there exists a point b of S and (taking b as origin) an orthonormal basis for L such that, with respect to this basis, S has equation \(\lambda_{1}\xi_{1}^{2} + \cdots + \lambda_{n-1}\xi_{n-1}^{2} + \xi_{n} = 0\) . (If \(\overline{S}\) is the projective quadric such that \(S = L \cap \overline{S}\) , c the pole (at infinity) with respect to \(\overline{S}\) of the hyperplane at infinity \(H_{0}\) , determine b by the condition that the line passing through b and c be perpendicular to the tangent hyperplane to S at the point b).

¶ 26) a) Let K be a commutative field, S a subset of K such that K is a maximal S-orderable field (chap.~VI, § 2, exerc. 8). Let f be a polynomial in K{[}X{]}, L its splitting field. Show that if, for every (total) order structure on K compatible with its ring structure, L admits an ordered extension structure of K, then necessarily L = K. (Argue by contradiction: proceeding as in exerc. 8 e) of chap.~VI, § 2, reduce to the case where one would have {[}L : K{]} = 2. Conclude by noting that if b ∈ K is not a square, there exists a total order structure on K, compatible with its ring structure, for which b \textless{} 0 (cf.~chap.~VI, § 2, n° 3, lemma of th. 1)).

\begin{enumerate}
\def\labelenumi{\alph{enumi})}
\setcounter{enumi}{1}
\tightlist
\item
  Extend the results of n° 3 (prop. 6 excepted) to the case where K is a maximal S-orderable field (*). (Begin by proving the analogue of prop. 5, using a). Then establish the analogue of cor. 2 of th. 2 for the case where the commutative algebra B consists of Hermitian endomorphisms, using prop. 10 of chap.~VIII, § 9, n° 4. Pass to the case of a normal endomorphism u for A = K(i) by noting that one can then write u = v + iw, where v and w are Hermitian and commute. Finally, if A is the quaternion field over K, E₀ the space E considered as a vector space of dimension 2n over K(i), and if one sets Φ(x, y) = Φ₁(x, y) + Φ₂(x, y)j, where Φ₁ and Φ₂ take their values in K(i), note that if u is normal for Φ, it is also normal for Φ₁.)
\end{enumerate}

¶ 27) Let K be a maximal ordered field, E a vector space of dimension n over K, Φ a non-degenerate symmetric bilinear form on E, of signature (s, t) distinct from (n, 0) and (0, n). Let (ei) be an orthogonal basis for Φ, such that Φ(ei, ei) = 1 for 1 ≤ i ≤ s, Φ(ei, ei) = -1 for s + 1 ≤ i ≤ n.~For every orthogonal transformation u ∈ U(Φ), let

\[
U = \left( \begin{array}{c c} M & N \\ P & Q \end{array} \right)
\]

the matrix of \(u\) with respect to \((e_{i})\) , written in the form of a square array of matrices corresponding to the partition of \([1, n]\) into \([1, s]\) and \([s + 1, n]\) .

\begin{enumerate}
\def\labelenumi{\alph{enumi})}
\tightlist
\item
  Prove the relations
\end{enumerate}

\[
\begin{array}{r l} ^ {t} M. M - ^ {t} P. P & = I _ {s} \\ ^ {t} Q. Q - ^ {t} N. N & = I _ {t} \\ ^ {t} M. N - ^ {t} P. Q & = 0. \end{array}
\]

\begin{enumerate}
\def\labelenumi{\alph{enumi})}
\setcounter{enumi}{1}
\item
  Let R be a matrix over K with t rows and s columns, such that \(I_{s} - ^{t}R \cdot R\) is the matrix of a positive non-degenerate Hermitian form (over \(K^{s}\) ). Show that det \((M + \lambda NR)\) does not change sign for \(-1 \leqslant \lambda \leqslant 1\) in K (show, using a), that \(^{t}(M + \lambda NR)(M + \lambda NR)\) is the matrix of a positive non-degenerate symmetric form).
\item
  Set \(\sigma(u) = \sigma(U) = \operatorname{sgn}(\det M)\) ; show that, for any two elements \(u\) , \(\nu\) of the orthogonal group \(\mathbf{U}(\Phi)\) , one has \(\sigma(u\nu) = \sigma(u)\sigma(\nu)\) (using \(a\) ) and \(b\) ). Deduce that the commutator subgroup of the special orthogonal group \(\mathbf{SU}(\Phi)\) is distinct from \(\mathbf{SU}(\Phi)\) (cf.~§ 10, exerc. 9).
\end{enumerate}

